\begin{center}[10.]\par\medskip\textit{Simplest Expressions for the Effects of Galvanic Currents.}\end{center}
The fundamental plane passes through the acting current element $AB$ and the point $C$ on which it acts.

Let the complex quantities denoting the positions $B,C$ relative to $A$ be $\beta,\gamma$, and let $r$ be the modulus of $\beta$. Finally, if there is already a current at $C$ with element $CD$, let $\gamma+\delta+i'\zeta$ be the complex quantity denoting the position of $D$ relative to $A$. Then one has

I.\quad if there is a current at $C$, for the force that $AB$ exerts on the material carrier of that current,
\[
\frac{\delta}{r}\quad\mathrm{Im.}\quad\frac{\beta}{\gamma}
\]

II.\quad if there is no current at $C$ but there is motion at $C$, denoted by $\varepsilon$ in the plane through $BC$, for the electromotive force,
\[
\frac{\beta}{r}\quad\mathrm R\quad\frac{\varepsilon}{\gamma}
\]
or in clearer notation
\SourceLeaf{621}\[
\begin{gathered}
G\quad\text{acting current element}\\[4pt]
\left.\begin{gathered}g\\m\end{gathered}\right\}\quad
\left.\begin{gathered}\text{existing current element}\\\text{existing motion}\end{gathered}\right\}\quad
\begin{gathered}\text{at the point}\\\text{acted upon}\end{gathered}\\[4pt]
\left.\begin{gathered}\gamma\\\mu\end{gathered}\right\}\quad
\text{force for producing}\Bigg\}\quad
\begin{gathered}\text{a current}\\\text{a motion}\end{gathered}\\[4pt]
r\quad\text{distance as the modulus of the complex quantity }l
\end{gathered}
\]
\[
\begin{aligned}
\text{1)}\qquad r\mu&=g\,.\,\mathrm{J.}\,\frac{G}{l}\\[4pt]
\text{2)}\qquad r\gamma&=G\,.\,\mathrm{R.}\,\frac{m}{l}
\end{aligned}
\]

\begin{center}[11.]\par\medskip\textit{Polygons with Straight Sides.}\end{center}
Let the point acted on be the origin. Then
\[
\text{I. For}\qquad X=\int\frac{\mathrm dx}{r}
\]
the contribution from the first side $PP'$ is
\[
\frac{x'-x}{r}\log\frac{\operatorname{cotg}\tfrac{1}{2}\theta'}{\operatorname{cotg}\tfrac{1}{2}\theta}
\]
where $\theta,\theta'$ are the angles between $PP'$ and $0P,0P'$, taking $PP'$ in the same direction. The quantity following the symbol log.\ is
\[
\begin{aligned}
&=\frac{r'+\dfrac{(x'-x)x'+(y'-y)y'+(z'-z)z'}{PP'}}{r+\dfrac{(x'-x)x+(y'-y)y+(z'-z)z}{PP'}}\\[6pt]
&=\frac{PP'\,.\,r'+r'r'-(r',PP')}{PP'\,.\,r-rr+(r,PP')}
\end{aligned}
\]

II.\quad For $V$, the contribution from the angle at $P'$ is the difference between the angle formed by the planes $0PP',0P'P''$ and $180^0$. If the area of a spherical triangle with sides $a,b,c$ is set $=\omega$, then
\[
\sin\tfrac{1}{2}\omega=\frac{6\text{ pyramid }0PP'P''}{4rr'r'\cos\tfrac{1}{2}a\,.\,\cos\tfrac{1}{2}b\,.\,\cos\tfrac{1}{2}c}
\]
\SourceLeaf{622}
\begin{center}[12.]\end{center}
\noindent$g$\quad Intensity of a galvanic current\par
\noindent$\mu$\quad Density of magnetism on a surface bounded by the rheophore\par
\noindent$\rho$\quad Distance of this surface from a second, negatively magnetized one
\[
\text{1.}\qquad g=\mu\rho
\]

2.\quad Action of one electric element on another. The position of the first relative to the second is specified by the complex quantity $u$ at time $t$, with distance $=r$.
\[
-\frac{u}{r^3}-\frac{\alpha}{ur}\left(\frac{\mathrm du}{\mathrm dt}\right)^2-\frac{\beta}{r}\frac{\mathrm d\mathrm du}{\mathrm dt^2}
\]

3.\quad Mutual action between an electric element and a magnetic element with relative velocity $v$
\[
\frac{\gamma v}{rr}
\]

4.\quad In a galvanic current of intensity $g$, during time $t$ the positive electricity $\varepsilon gt$ passes through every cross-section in one direction, and the negative electricity $-\varepsilon gt$ in the other.

5.\quad The aim is to determine the relations between $\alpha,\beta,\gamma,\varepsilon$. The relation $\gamma=2\beta\varepsilon$ follows from induction as a current forms over the time $\theta$, and for as long as the current is forming, the force becomes —
\[
\begin{gathered}
=\frac{\gamma\pi hhg}{\theta RR}=\frac{2\beta\pi hh\varepsilon g}{\theta RR}\\[6pt]
\gamma=4\alpha\varepsilon,\qquad\frac{\gamma\pi hhgv}{R^3}=\frac{4\alpha\pi hh\varepsilon gv}{R^3}\\[6pt]
2\gamma\varepsilon=1,\qquad\alpha=\frac{1}{8\varepsilon\varepsilon}=\tfrac{1}{2}\gamma\gamma,\qquad\beta=\frac{1}{4\varepsilon\varepsilon}=\gamma\gamma
\end{gathered}
\]

